\documentclass[10pt,a4paper]{amsart}
\usepackage[margin=19mm]{geometry}
\usepackage{amsmath,amssymb,mathtools}
\usepackage[scr=rsfs,cal=euler]{mathalfa}
\usepackage{xfrac}
\usepackage[unicode,hidelinks,bookmarks=false]{hyperref}
\newcommand{\ie}{i.e.\ }
\newcommand{\cf}{cf.\ }
\newcommand{\resp}{respectively\ }
\newcommand{\Subsection}{\subsection}
\providecommand{\nobreakdash}{\nobreak\hbox{-}}
\setlength{\emergencystretch}{2em}
\multlinegap0pt
\usepackage{bm}
\usepackage{bbm}
\usepackage{dsfont}
\usepackage{xspace}
\usepackage{cancel}
\usepackage{mathtools}
\usepackage{amsthm}
\makeatletter
\@ifundefined{lemma}{\newtheorem{lemma}{Lemma}}{}
\makeatother
\makeatletter
\newcommand{\E}{\mathbb{E}}
\newcommand{\Var}{\mathbb{V}ar}
\expandafter\ifx\csname customex\endcsname\relax
\newenvironment{customex}[1]
  {\renewcommand\theinnercustomex{#1}\innercustomex}
  {\endinnercustomex}
\fi
\makeatother
\title[On the Lower Confidence Band for the Optimal Welfare in Poli]{On the Lower Confidence Band for the Optimal Welfare in Policy Learning}
\author{Kirill Ponomarev, Vira Semenova}
\date{Source: arXiv:2410.07443v3 (CC BY 4.0)}
\begin{document}
\maketitle

\noindent\emph{Source and licence.} On the Lower Confidence Band for the Optimal Welfare in Policy Learning. Version arXiv:2410.07443v3 (CC BY 4.0), distributed under the Creative Commons Attribution 4.0 International licence, as recorded in the arXiv licence metadata for this exact version and shown on the abstract page https://arxiv.org/abs/2410.07443v3. Full rights notice: licenses/arxiv-2410.07443-CC-BY-4.0.txt.\par\smallskip

\noindent\textit{Editorial scope.} Counted unit: the Lemma whose source label is \texttt{lem:bias} in the article (source0.tex lines 1086--1093, source label \texttt{lem:bias}), delivered together with the article's own complete original proof of it (source0.tex lines 1095--1151, one \texttt{proof} environment of 57 lines). No mathematics is written, repaired or re-derived: statement and proof are the article's own text line for line. Required context retained uncounted: eq:mean1 (lines 1068--1071); eq:mean3 (lines 1068--1071); eq:chernoff (lines 1073--1080). Every definition, display and label of those lines is retained unchanged. Omitted from this package and not redistributed: 1--1067, 1072--1072, 1081--1085, 1094--1094, 1152--1545. Nothing inside a retained excerpt is omitted or altered: every retained body is delivered exactly as the article's lines are written, so no omission receipt of any kind accompanies this package. Citations: the retained excerpts carry no citation command at all, so no bibliography entry of the article is transcribed and no reference is redistributed. Reference adaptation: none was needed and none was made; every reference inside the delivered unit resolves inside it, so no unresolved reference or citation is produced. Printed slips kept literal: no printed word, symbol, hypothesis or number of the counted statement or of its proof was altered. Layout and class: the article's own class is \texttt{article} with packages including babel, amsmath, amsfonts, amssymb, amsthm, mathtools, bm, bbm, dsfont, xcolor, colortbl, chngcntr, apptools, comment; the wrapper is the lane's pinned \texttt{amsart} editorial wrapper at 10pt on A4, which restates the theorem-like environments the retained text needs and adds the article's own macro definitions that the retained text needs (\texttt{\textbackslash E}, \texttt{\textbackslash Var}), with extra wrapper packages (bm, bbm, dsfont, xspace, cancel, mathtools). The delivered numbering is the unit's own. Assets: no retained span contains a figure environment, a graphics-inclusion command, a PDF-page inclusion, a table or a TikZ picture; no image file is copied into this package, the package declares no asset dependency and no image file is redistributed with it. Licence: version arXiv:2410.07443v3 of this article is distributed under the Creative Commons Attribution 4.0 International licence, as recorded in the arXiv licence metadata for this exact version and shown on the abstract page https://arxiv.org/abs/2410.07443v3. A full rights notice is delivered at licenses/arxiv-2410.07443-CC-BY-4.0.txt. Characters: every retained excerpt is pure ASCII, so the delivered engine, XeLaTeX with its default Latin Modern text font, reports no missing character.
\begin{align}
\E [  (N_{dx} + 1)^{-1}   ] & = \dfrac{4}{(N+1) } - \dfrac{4}{(N+1)  } (3/4)^{N+1}  \leq 4/N \label{eq:mean1} \\
 \E [   (N_{dx} + 1)^{-2}   ]   & \leq 32  N^{-2}. \label{eq:mean3}
\end{align} 

\begin{equation} \label{eq:chernoff}
\begin{array}{cl}
	P\left(\left|\hat{p} - 1/2\right| \geq C_N\right) &= P(\hat{p} \geq 1/2 + C_N) + P(\hat{p} \leq 1/2 - C_N)\\[3mm]
	&= P(S \geq (1 + 2C_N)\frac{N}{2}) + P(S \leq (1-2C_N)\frac{N}{2})\\[3mm]
	&\leq \exp \left(-\frac{2}{3}N(C_N)^2 \right) + \exp\left(-N(C_N)^2 \right)\\[3mm]
	&\leq 2\exp \left(-\frac{2}{3}N(C_N)^2 \right).
\end{array}
\end{equation}

 \begin{lemma}
   \label{lem:bias}
 For $N \geq 100$ and $C_N = \sqrt{2.25 \ln N/N}$, the following bounds hold for any $d,x \in \{1,0\}$
\begin{align}
 (1-10 C_N) < \E [\sigma^{-2}(d,1)  N \cdot \Var ( \widehat m_{d1} \mid \mathbf{X}, \mathbf{D})\cdot    \widehat{p}^2] <  (1+ 10 C_N). \label{eq:varbound} \\
(1-10 C_N) <   \E [\sigma^{-2}(d,0) N \cdot \Var ( \widehat m_{d0} \mid \mathbf{X}, \mathbf{D})   \cdot (1-\widehat{p})^2] < (1+ 10 C_N). \label{eq:varbound2} 
\end{align}
\end{lemma}

\begin{proof}

\textbf{Step 1 (Notation).} Denote the expression inside the expectation of \eqref{eq:varbound} by
$$
\Xi_d = \sigma^{-2}(d,1)  N \Var ( \widehat m_{d1} \mid \mathbf{X}, \mathbf{D})   \widehat{p}^2 = N N_{d1}(N_{d1}  + 1)^{-2} \widehat{p}^2,
$$
and note that its probability limit  as $N \to \infty$ equals $1$. Denoting 
$$\psi^1_{d} (t)= N t  (N_{d1} + 1)^{-1}; \;\; \quad \psi^2_{d} (t)= N t  (N_{d1} + 1)^{-2},
$$ 
we can decompose the asymptotic error as 
\begin{align}
\label{eq:vardecomp}
\Xi_d -1 &= N N_{d1}(N_{d1}  + 1)^{-2} \widehat{p}^2-1\\[2mm]
&= \psi^1_{d} (\widehat{p}^2) - \psi^2_{d}(\widehat{p}^2) - 1 \nonumber \\[2mm]
&=  \psi^1_{d} (\widehat{p}^2 -p^2) + \psi^1_d (p^2) -  \psi^2_{d}(\widehat{p}^2) - 1 \nonumber \\[2mm]
&=\underbracket{\psi^1_{d} (\widehat{p}^2-p^2)}_{S_2}   + \underbracket{\psi^1_{d} (p^2) -  N/(N+1) }_{S_1} - \underbracket{\psi^2_{d}(\widehat{p}^2)}_{S_3}  -\underbracket{1/(N+1)}_{S_4}. \nonumber 
\end{align}

\noindent 
\textbf{Step 2 (Leading term $S_2$).}  Recall that $p = 1/2$. On the event $\mathcal{M}_N = \{ |\widehat p - 1/2| <  C_N \}$, the error $ |\widehat p^2 -1/4| \leq |\widehat{p} - 1/2||\widehat{p} + 1/2| \leq 1.5C_N$. As a result, 
\begin{align}
  |\E\left[S_2 \bm{1}\{ \mathcal{M}_N  \}\right] | &\leq  \E \left[ \psi^1_{d} (|\widehat p^2 - p^2|)  \bm{1}\{ \mathcal{M}_N  \}\right]  \nonumber \\[2mm]
   &\leq \E [ \psi^1_{d} (1.5C_N)   \bm{1}\{ \mathcal{M}_N  \}] \nonumber  \\[2mm]
   &\leq 1.5C_N \E [ \psi_{d}^1(1)] \nonumber \\[2mm]
   & \leq 6C_N, \label{eq:6cn}
\end{align}
 where the first three lines follow from linearity of $\psi^1_d(\cdot)$ and monotonicity of expectation and the last one follows from \eqref{eq:mean1}.  On the  event $\mathcal{M}_N^c$, we can bound $ |\widehat p^2 -p^2| \leq 1 \text{, a.s.,}$ so that
\begin{align*}
    |\E [S_2 \bm{1}\{ \mathcal{M}^c_N  \}] | &\leq \E [ \psi^1_{d} (|\widehat p^2 - p^2|)  \bm{1}\{ \mathcal{M}^c_N  \}]  \\[2mm]
   & \leq \E [ \psi^1_{d} (1)  \bm{1}\{ \mathcal{M}^c_N  \}] \\[2mm]
   & \leq N  P (\mathcal{M}_N^c), 
   \end{align*}   
   where the last line follows from $N_{d1}\geq 0$ and $(N_{d1} +1)^{-1} \leq 1$, a.s.. Using \eqref{eq:chernoff} and $C_N = \sqrt{ 2.25 \ln N/N}$,
\begin{align}
 \textstyle N P (\mathcal{M}^c_N)  \leq 2 N \exp(-\frac{2}{3}N (C_N)^2 ) = 2 N^{-1/2} \leq C_N, \qquad \forall N \geq 6. \label{eq:7cn}
\end{align}
Adding \eqref{eq:6cn} and \eqref{eq:7cn} gives
$ |\E[S_2]|  \leq 7 C_N. $ 

\medskip  

\noindent \textbf{Step 3 (Terms $S_1, S_3, S_4$).} 
 Note that $S_4 =(N+1)^{-1} \leq C_N$.   Invoking \eqref{eq:mean1} gives
  $$
   |\E [S_1] | = 1/4 |\E [ \psi^1_{d} (1) - 4N/(N+1)] | =  N/(N+1) (3/4)^{N+1} \leq   C_N, \quad \forall N \geq 2.
   $$ 
Invoking \eqref{eq:mean3} gives
 $$
  0 \leq \E [S_3]  =\E [ \psi^2_{d} (\widehat{p}^2) ] \leq \E [ \psi^2_{d} (1) ]  \leq  32   N^{-1} \leq  C_N, \quad \forall N \geq 100.
  $$ 
Combining the bounds gives
$$
|\E [ \Xi_d -1] | \leq \sum_{j=1}^4 |\E [ S_d ]| \leq 10 C_N.
$$

\noindent \textbf{Step 4 (Conclusion).} Steps 1--3 established \eqref{eq:varbound}, which corresponds to $x=1$.  The symmetry of DGPs implies \eqref{eq:varbound2} with $x=0$. 
\end{proof}

\end{document}
